\section*{T?M T?T B?O C?O}
\phantomsection\addcontentsline{toc}{section}{\numberline {}T?M T?T B?O C?O}

B?o c?o n?y tr?nh b?y m?t h? th?ng ph?t hi?n gian l?n th? t?n d?ng s? d?ng m? h?nh XGBoost tr?n d? li?u giao d?ch c? m?c m?t c?n b?ng cao. D? li?u ???c chia th?nh 16 chunk theo th? t? th?i gian; chunk 01 ???c s? d?ng ?? ph?n t?ch EDA, l?a ch?n ??c tr?ng v? hu?n luy?n ban ??u, trong khi c?c chunk 02--16 ???c d?ng ?? ??nh gi? kh? n?ng t?ng qu?t h?a v? ph?t hi?n drift.

K?t qu? ph?n t?ch d? li?u cho th?y t? l? gian l?n trong chunk 01 ch? chi?m 1.0597\%, do ?? accuracy kh?ng ph? h?p l?m th??c ?o ch?nh. C?c ??c tr?ng quan tr?ng nh?t bao g?m s? ti?n giao d?ch, category, merchant frequency, job frequency v? m?t s? ??c tr?ng h?nh vi theo th?i gian. Qu? tr?nh l?a ch?n ??c tr?ng gi?m s? bi?n ??u v?o t? 41 xu?ng 25 nh?ng ch? l?m gi?m PR-AUC 0.0013, cho th?y nhi?u ??c tr?ng ban ??u l? d? th?a.

Sau khi tuning, hai m? h?nh XGBoost t?t nh?t ??t PR-AUC tr?n 0.96 v? F1-score tr?n 0.95 tr?n test set c?a chunk 01. Tuy nhi?n, khi ??nh gi? tr?n c?c chunk sau, m?t s? chunk c? F1-score th?p. Ph?n t?ch chi ti?t cho th?y chunk 5--6 c? d?u hi?u concept ho?c label drift, trong khi chunk 10--11 ch? y?u g?p threshold calibration drift. V? v?y, b?o c?o ?? xu?t s? d?ng \texttt{xgb\_class\_weight} v?i selected features, dynamic threshold v? walk-forward validation cho tri?n khai th?c t?.

\newpage

\section*{THESIS SUMMARY}
\phantomsection\addcontentsline{toc}{section}{\numberline {}THESIS SUMMARY}

This report presents a credit card fraud detection system based on XGBoost for highly imbalanced tabular transaction data. The dataset is split into 16 time-sorted chunks. Chunk 01 is used for exploratory analysis, feature selection and initial training, while chunks 02--16 are used to evaluate temporal generalization and drift.

The EDA shows that fraud accounts for only 1.0597\% of transactions in chunk 01, making accuracy unsuitable as the main metric. Transaction amount and category are the strongest signals, while geographic distance and time-of-day are weak standalone predictors. Feature selection reduces the model-ready feature set from 41 to 25 features with almost no PR-AUC loss.

Tuned XGBoost models perform strongly on chunk 01, but later chunks reveal temporal instability. Chunks 5--6 are likely concept or label drift candidates, while chunks 10--11 mainly show threshold calibration drift. The recommended pipeline is XGBoost with class weighting, selected features, dynamic thresholding, walk-forward validation and drift-aware monitoring.
